\clearpage
\section{Problemas Resueltos del Tema 1}

\begin{question}{Módulo de elasticidad a partir de alargamiento en barra recta}{prob:1}
    Una barra recta de sección uniforme está sometida a una fuerza de tracción axial. La sección transversal es de $6\text{ cm}^2$ y la longitud inicial es de $4\text{ m}$. Si el alargamiento total registrado bajo una carga axial de $12\,600\text{ kgf}$ es de $0.40\text{ cm}$, determinar el módulo de elasticidad del material en $\text{kgf}/\text{cm}^2$.
\end{question}

\solve{}
\textbf{1. Datos del problema:}
\begin{itemize}
    \item Sección transversal: $A_0 = 6\text{ cm}^2$
    \item Longitud inicial: $L_0 = 4\text{ m} = 400\text{ cm}$
    \item Carga axial aplicada: $P = 12\,600\text{ kgf}$
    \item Alargamiento elástico total: $\delta = 0.40\text{ cm}$
\end{itemize}

\textbf{2. Ecuaciones fundamentales:}
Asumiendo comportamiento elástico lineal bajo la Ley de Hooke:
\begin{equation*}
    \sigma = E \, \varepsilon \quad \Longleftrightarrow \quad \frac{P}{A_0} = E \, \frac{\delta}{L_0}
\end{equation*}
Despejando el módulo de Young $E$:
\begin{equation*}
    E = \frac{P \cdot L_0}{A_0 \cdot \delta}
\end{equation*}

\textbf{3. Sustitución numérica:}
\begin{equation*}
    E = \frac{12\,600\text{ kgf} \times 400\text{ cm}}{6\text{ cm}^2 \times 0.40\text{ cm}} = \frac{5\,040\,000}{2.4} = 2.10 \times 10^6\text{ kgf}/\text{cm}^2
\end{equation*}
En unidades del Sistema Internacional:
\begin{equation*}
    E \approx 2.10 \times 10^6 \times 9.80665 \times 10^4\text{ Pa} \approx 206\text{ GPa} \quad (\text{correspondiente a un acero})
\end{equation*}


\vspace{1em}

\begin{question}{Compresión uniaxial en cilindro hueco}{prob:2}
    Un cilindro hueco recto de sección circular tiene un diámetro exterior de $7.5\text{ cm}$ y un diámetro interior de $6.0\text{ cm}$. Si se le aplica una fuerza axial de compresión de $5\,000\text{ kgf}$, determinar:
    \begin{enumerate}[label=\alph*)]
        \item La tensión normal de compresión desarrollada.
        \item El acortamiento total que experimenta en una longitud inicial de $50\text{ cm}$.
    \end{enumerate}
    Considérese un módulo de elasticidad $E = 1.05 \times 10^6\text{ kgf}/\text{cm}^2$ y despréciese toda probabilidad de pandeo lateral.
\end{question}

\solve{}
\textbf{1. Datos geométricos y mecánicos:}
\begin{itemize}
    \item Diámetro exterior: $D_{\text{ext}} = 7.5\text{ cm}$
    \item Diámetro interior: $D_{\text{int}} = 6.0\text{ cm}$
    \item Fuerza axial de compresión: $P = 5\,000\text{ kgf}$
    \item Longitud inicial: $L_0 = 50\text{ cm}$
    \item Módulo elástico: $E = 1.05 \times 10^6\text{ kgf}/\text{cm}^2$
\end{itemize}

\textbf{2. Cálculo de la sección transversal neta:}
\begin{equation*}
    A_0 = \frac{\pi}{4} \left(D_{\text{ext}}^2 - D_{\text{int}}^2\right) = \frac{\pi}{4} \left(7.5^2 - 6.0^2\right) = \frac{\pi}{4} (56.25 - 36.00) = \frac{\pi}{4} (20.25) \approx 15.904\text{ cm}^2
\end{equation*}

\textbf{3. Tensión normal de compresión ($\sigma$):}
\begin{equation*}
    \sigma = \frac{P}{A_0} = \frac{5\,000\text{ kgf}}{15.904\text{ cm}^2} \approx 314.38\text{ kgf}/\text{cm}^2 \approx 314.4\text{ kgf}/\text{cm}^2
\end{equation*}

\textbf{4. Acortamiento total elástico ($\delta$):}
Aplicando la relación de Hooke $\sigma = E \, (\delta / L_0)$:
\begin{equation*}
    \delta = \frac{\sigma \cdot L_0}{E} = \frac{314.38\text{ kgf}/\text{cm}^2 \times 50\text{ cm}}{1.05 \times 10^6\text{ kgf}/\text{cm}^2} \approx 0.01497\text{ cm} \approx 0.0150\text{ cm} = 0.15\text{ mm}
\end{equation*}


\vspace{1em}

\begin{question}{Deformación transversal elástica y coeficiente de Poisson en latón}{prob:3}
    Se aplica una carga uniaxial de tracción en la dirección longitudinal de una barra cilíndrica de latón que posee un módulo elástico $E = 10.1 \times 10^4\text{ MPa}$, coeficiente de Poisson $\nu = 0.35$ y un diámetro inicial $d_0 = 10\text{ mm}$. Determinar la magnitud de la carga $P$ necesaria para producir una reducción en el diámetro de $2.5 \times 10^{-3}\text{ mm}$, suponiendo deformación puramente elástica.
\end{question}

\solve{}
\textbf{1. Identificación de magnitudes y datos:}
\begin{itemize}
    \item Módulo de elasticidad: $E = 10.1 \times 10^4\text{ MPa} = 101\,000\text{ N}/\text{mm}^2$
    \item Coeficiente de Poisson: $\nu = 0.35$
    \item Diámetro inicial: $d_0 = 10\text{ mm}$
    \item Variación transversal del diámetro: $\Delta d = -2.5 \times 10^{-3}\text{ mm}$
\end{itemize}

\textbf{2. Deformación unitaria transversal y longitudinal:}
La deformación unitaria lateral en el diámetro es:
\begin{equation*}
    \varepsilon_x = \frac{\Delta d}{d_0} = \frac{-2.5 \times 10^{-3}\text{ mm}}{10\text{ mm}} = -2.5 \times 10^{-4}
\end{equation*}
A partir del coeficiente de Poisson ($\nu = -\varepsilon_x / \varepsilon_z$):
\begin{equation*}
    \varepsilon_z = -\frac{\varepsilon_x}{\nu} = -\frac{-2.5 \times 10^{-4}}{0.35} \approx 7.1429 \times 10^{-4}
\end{equation*}

\textbf{3. Tensión axial requerida ($\sigma_z$):}
Por la ley de Hooke en régimen elástico:
\begin{equation*}
    \sigma_z = E \cdot \varepsilon_z = 101\,000\text{ MPa} \times 7.1429 \times 10^{-4} \approx 72.143\text{ MPa} = 72.143\text{ N}/\text{mm}^2
\end{equation*}

\textbf{4. Cálculo de la carga axial $P$:}
La sección transversal inicial de la probeta es:
\begin{equation*}
    A_0 = \frac{\pi}{4} d_0^2 = \frac{\pi}{4} (10\text{ mm})^2 = 25\pi\text{ mm}^2 \approx 78.54\text{ mm}^2
\end{equation*}
Por tanto, la carga $P$ aplicada resulta:
\begin{equation*}
    P = \sigma_z \cdot A_0 = 72.143\text{ N}/\text{mm}^2 \times 78.54\text{ mm}^2 \approx 5\,666\text{ N} \quad (5.67\text{ kN})
\end{equation*}
\textit{Nota sobre el resultado:} Si el cambio diametral fuese $\Delta d = 2.5 \times 10^{-5}\text{ mm}$, la carga sería $56.66\text{ N}$. Para el valor enunciado en el texto de $2.5 \times 10^{-3}\text{ mm}$, la carga elástica rigurosa es $P = 5\,666\text{ N}$.


\vspace{1em}

\begin{question}{Ductilidad y tensión real de fractura en probeta de acero}{prob:4}
    Una probeta cilíndrica de acero con un diámetro original $d_0 = 12.8\text{ mm}$ se ensaya a tracción uniaxial hasta la fractura, alcanzando una resistencia nominal a la rotura $\sigma_D = 460\text{ MPa}$. Una vez rota la probeta, se mide el diámetro transversal mínimo en la sección de fractura resultando $d_f = 10.7\text{ mm}$. Evaluar:
    \begin{enumerate}[label=\alph*)]
        \item La ductilidad del material expresada mediante el porcentaje de estricción (\%ES).
        \item La tensión real $\sigma_{\text{real}}$ que soportaba la sección de estricción en el instante exacto de la fractura.
    \end{enumerate}
\end{question}

\solve{}
\textbf{1. Datos iniciales y finales:}
\begin{itemize}
    \item Diámetro inicial: $d_0 = 12.8\text{ mm}$
    \item Diámetro final en rotura: $d_f = 10.7\text{ mm}$
    \item Resistencia nominal a rotura: $\sigma_D = 460\text{ MPa} = 460\text{ N}/\text{mm}^2$
\end{itemize}

\textbf{2. Cálculo de áreas transversales inicial y final:}
\begin{align*}
    A_0 & = \frac{\pi}{4} d_0^2 = \frac{\pi}{4} (12.8)^2 \approx 128.68\text{ mm}^2 \\
    A_f & = \frac{\pi}{4} d_f^2 = \frac{\pi}{4} (10.7)^2 \approx 89.92\text{ mm}^2
\end{align*}

\textbf{3. Evaluación del porcentaje de estricción (\%ES):}
\begin{equation*}
    \%ES = \frac{A_0 - A_f}{A_0} \times 100 = \frac{128.68 - 89.92}{128.68} \times 100 = \frac{38.76}{128.68} \times 100 \approx 30.1\% \approx 30\%
\end{equation*}

\textbf{4. Cálculo de la tensión real de fractura ($\sigma_{\text{real}}$):}
La carga total $P_D$ actuante en el momento del colapso se deduce a partir de la tensión nominal de rotura referida al área original:
\begin{equation*}
    P_D = \sigma_D \cdot A_0 = 460\text{ N}/\text{mm}^2 \times 128.68\text{ mm}^2 \approx 59\,192.8\text{ N}
\end{equation*}
La tensión real en la rotura se evalúa dividiendo dicha carga entre la sección transversal final real contraída $A_f$:
\begin{equation*}
    \sigma_{\text{real}} = \frac{P_D}{A_f} = \frac{\sigma_D \cdot A_0}{A_f} = 460 \times \frac{128.68}{89.92} = 460 \times \left(\frac{12.8}{10.7}\right)^2 \approx 460 \times 1.431 \approx 658.2\text{ MPa}
\end{equation*}

\clearpage
\section{Test de Autoevaluación - Tema 1}

\begin{enumerate}[label=\textbf{\arabic*.}]
    \item Los materiales compuestos .....
          \begin{enumerate}[label=\alph*.]
              \item Tienen propiedades homogéneas en todas las direcciones del espacio
              \item \textcolor{red!85!black}{Muestran un buen comportamiento frente a cargas cíclicas.}
              \item Presentan una resistencia específica inferior a la de los metales.
          \end{enumerate}

    \item Si se desea el uso de un material que requiera tener baja densidad y alta flexibilidad, de forma general recurriríamos al uso de un material....
          \begin{enumerate}[label=\alph*.]
              \item Cerámico
              \item \textcolor{red!85!black}{Polimérico termoplástico.}
              \item Polimérico termoestable
          \end{enumerate}

    \item Cuando un material sufre una deformación elástica...
          \begin{enumerate}[label=\alph*.]
              \item \textcolor{red!85!black}{Recupera sus dimensiones al dejar de aplicarse la carga.}
              \item Recupera sus dimensiones siempre que no se haya superado el límite de fluencia o cedencia.
              \item Aumenta su límite de proporcionalidad tras dejar de aplicarse la carga.
          \end{enumerate}

    \item Señale la respuesta correcta:
          \begin{enumerate}[label=\alph*.]
              \item Un mayor valor del módulo de resiliencia indica que el material tiene mayor resistencia a la deformación plástica.
              \item \textcolor{red!85!black}{Un mayor valor del Alargamiento a la rotura y de la Estricción indica que el material es más dúctil.}
              \item Ambas respuestas son correctas
          \end{enumerate}

    \item Señale la respuesta correcta:
          \begin{enumerate}[label=\alph*.]
              \item \textcolor{red!85!black}{El objetivo del ensayo de tracción normalizado permite comparar el comportamiento de materiales entre sí.}
              \item El ensayo tracción-deformación suele mostrar la tensión real que sufre un material en función de su deformación unitaria.
              \item El objetivo del ensayo de tracción normalizado es evaluar el comportamiento que tendrá un material en su uso real.
          \end{enumerate}

    \item Señale la respuesta correcta:
          \begin{enumerate}[label=\alph*.]
              \item \textcolor{red!85!black}{La cianuración es un tratamiento superficial que pretende aumentar la resistencia mecánica de la superficie del acero mediante la inclusión de grupos $\text{CN}^-$ (cianuros) sobre la superficie.}
              \item La nitruración es un tratamiento superficial que pretende aumentar la resistencia mecánica de la superficie del acero mediante la formación de nitratos con base férrea.
              \item La cementación es un tratamiento superficial que pretende aumentar la resistencia mecánica de la superficie del acero mediante la formación de carburo de hierro.
          \end{enumerate}

    \item Señale la respuesta correcta:
          \begin{enumerate}[label=\alph*.]
              \item \textcolor{red!85!black}{Ambas respuestas son correctas}
              \item Un aumento en el contenido en carbono de los aceros los hace menos tenaces.
              \item Un material es más tenaz cuando es capaz de absorber mayor energía plástica y elástica antes de su rotura.
          \end{enumerate}

    \item Los materiales compuestos .....
          \begin{enumerate}[label=\alph*.]
              \item Están formados por un material que le confiere rigidez (fibras) y otro que actúa de relleno.
              \item \textcolor{red!85!black}{Suelen tener una resistencia mecánica intermedia entre las fibras y la matriz}
              \item Ambas respuestas son correctas
          \end{enumerate}

    \item Señale la respuesta correcta:
          \begin{enumerate}[label=\alph*.]
              \item Ambas respuestas son correctas
              \item La tenacidad es la capacidad de un material de absorber energía plástica antes de su fractura.
              \item \textcolor{red!85!black}{La tenacidad se puede evaluar mediante ensayos que miden la energía de impacto de un martillo sobre el material}
          \end{enumerate}

    \item Señale la respuesta correcta:
          \begin{enumerate}[label=\alph*.]
              \item \textcolor{red!85!black}{La deformación plástica aumenta el límite elástico de un material}
              \item La deformación plástica disminuye el límite de proporcionalidad de un material.
              \item La deformación plástica reduce el punto de fluencia de un material
          \end{enumerate}

    \item Señale la respuesta correcta:
          \begin{enumerate}[label=\alph*.]
              \item Un aumento de la carga en el ensayo de termofluencia conlleva una reducción de la mínima rapidez de termofluencia
              \item \textcolor{red!85!black}{La termofluencia conlleva el endurecimiento del material debido a la deformación plástica sufrida.}
              \item La mínima rapidez de termofluencia se produce en las etapas iniciales del ensayo de termofluencia, y nos permitirá evaluar el tiempo de vida que tendría un equipo.
          \end{enumerate}

    \item Señale la respuesta correcta:
          \begin{enumerate}[label=\alph*.]
              \item \textcolor{red!85!black}{El ensayo de Erichsen permite evaluar la capacidad que tiene un material plano de ser transformado en piezas huecas como fondos, carcasas, etc.}
              \item El ensayo de flexibilidad permite evaluar las propiedades mecánicas de resistencia que tiene un material.
              \item Dentro de los ensayos no destructivos podemos encontrar el ensayo de dureza Rockwell superficial
          \end{enumerate}

    \item Señale la respuesta correcta:
          \begin{enumerate}[label=\alph*.]
              \item El límite de proporcionalidad y el límite de elasticidad siempre coinciden.
              \item \textcolor{red!85!black}{El límite de fluencia inferior es un punto característico del material}
              \item Ambas respuestas son correctas
          \end{enumerate}

    \item En cuanto a la resistencia por fatiga que tienen los materiales
          \begin{enumerate}[label=\alph*.]
              \item La resistencia a la fatiga dependerá de la resistencia mecánica propia que tenga el material, pero no del ambiente de trabajo en el que desarrolle su labor.
              \item \textcolor{red!85!black}{Superficies más pulidas conllevarán mayor resistencia a la fatiga.}
              \item Los procesos superficiales de carburación, nitruración, etc. no mejora la resistencia a la fatiga del material ya que la resistencia mecánica del material completo sigue siendo la misma.
          \end{enumerate}

    \item La tensión ingenieril se define como:
          \begin{enumerate}[label=\alph*.]
              \item Ambas son incorrectas
              \item \textcolor{red!85!black}{Cociente entre la carga aplicada y la sección (inicial) del material transversal a la dirección de aplicación de la carga.}
              \item Cociente entre la carga aplicada y la sección (en cada momento del ensayo) del material transversal a la dirección de aplicación de la carga.
          \end{enumerate}

    \item Señale la respuesta correcta:
          \begin{enumerate}[label=\alph*.]
              \item Un aumento de la temperatura hace que generalmente los materiales sean más rígidos.
              \item \textcolor{red!85!black}{Los ensayos de termofluencia permiten evaluar el comportamiento de un material frente a temperaturas y cargas elevadas en función del tiempo.}
              \item Ambas respuestas son correctas
          \end{enumerate}

    \item Señale la respuesta correcta:
          \begin{enumerate}[label=\alph*.]
              \item El laminado del acero tiene como objetivo la reducción del espesor de la plancha de metal
              \item \textcolor{red!85!black}{El carácter inoxidable de un acero se lo da el contenido en cromo en la aleación, que suele estar entre en una proporción de 10-27\%.}
              \item Uno de los aceros inoxidables más utilizados en la industria alimentaria es el AISI 9228.
          \end{enumerate}

    \item Señale la respuesta correcta:
          \begin{enumerate}[label=\alph*.]
              \item La deformación unitaria se define como el cociente entre el alargamiento y la longitud inicial de la probeta
              \item \textcolor{red!85!black}{Ambas son correctas}
              \item El alargamiento o deformación se define como la elongación que sufre un material en la dirección de aplicación de la carga
          \end{enumerate}

    \item Señale la respuesta correcta:
          \begin{enumerate}[label=\alph*.]
              \item La resistencia máxima a la tracción es el punto donde la tensión real aplicada es máxima
              \item Ambas respuestas son correctas
              \item \textcolor{red!85!black}{En caso de que el límite de fluencia no esté claro, se puede definir como la tensión para la cual la deformación plástica es del 0.2\%.}
          \end{enumerate}

    \item Señale la respuesta correcta:
          \begin{enumerate}[label=\alph*.]
              \item Un aumento de la temperatura en el ensayo de termofluencia conlleva una reducción de la mínima rapidez de termofluencia
              \item La fluencia lenta o termofluencia es mayor en las etapas iniciales del ensayo
              \item \textcolor{red!85!black}{La mínima rapidez de termofluencia nos muestra el crecimiento exponencial de la deformación que sufre un material.}
          \end{enumerate}

    \item Si se desea el uso de un material que requiera tener alta dureza, resistencia al desgaste y con propiedades aislantes y térmicas, de forma general recurriríamos al uso de un material....
          \begin{enumerate}[label=\alph*.]
              \item Polimérico
              \item Metálico
              \item \textcolor{red!85!black}{Cerámico}
          \end{enumerate}

    \item Señale la respuesta correcta:
          \begin{enumerate}[label=\alph*.]
              \item La tensión de trabajo coincide con la tensión en el punto de fluencia o con la resistencia máxima a la tracción, seleccionándose uno de estos valores en función de las condiciones de servicio del material.
              \item \textcolor{red!85!black}{La tensión de seguridad o tensión de trabajo es la tensión máxima que puede soportar un material en sus condiciones de servicio.}
              \item El producto de la tensión en el punto de fluencia o en el punto de máxima tracción (dependiendo del material y uso del mismo) por el coeficiente de seguridad nos determinará la tensión de trabajo del material.
          \end{enumerate}

    \item El acero es una aleación hierro-carbono donde.....
          \begin{enumerate}[label=\alph*.]
              \item Cuanto mayor porcentaje en carbono mayor presente, mayor resistencia mecánica tendrá el material
              \item \textcolor{red!85!black}{Ambas respuestas son correctas}
              \item La ferrita es hierro casi puro y la cementita es un carburo de hierro
          \end{enumerate}

    \item Señale la respuesta correcta:
          \begin{enumerate}[label=\alph*.]
              \item Ambas respuestas son correctas
              \item \textcolor{red!85!black}{La ductilidad es una medida de la deformación plástica que sufre un material antes de su fractura.}
              \item La tenacidad es una medida de la deformación elástica que sufre un material antes de su fractura.
          \end{enumerate}

    \item Señale la respuesta correcta:
          \begin{enumerate}[label=\alph*.]
              \item \textcolor{red!85!black}{Los tratamientos mecánicos tienen como objetivo producir una deformación plástica en el material para conferir una resistencia determinada.}
              \item Un enfriamiento rápido del material desde la temperatura crítica máxima tiene como objetivo reducir las tensiones residuales que aparecen en el material tras un proceso térmico más severo.
              \item Los tratamientos superficiales aplicados a aceros pretenden dar al material una mayor resistencia mecánica al material debido a la deformación plástica que sufre la superficie.
          \end{enumerate}

    \item Ensayo de microdureza:
          \begin{enumerate}[label=\alph*.]
              \item \textcolor{red!85!black}{Vickers y Knoop}
          \end{enumerate}
\end{enumerate}

